Le sel de \textsc{Mohr} est un composé stable, peu oxydable, utilisé pour
préparer des solutions étalons. Il se présente sous forme de cristaux verdâtres.
Sur l'étiquette d'un flacon de sel de \textsc{Mohr}, on peut lire~:

\begin{itemize}
    \item ammonium-fer(II) sulfate à 6 molécules d'eau~;
    \item \ce{H8FeN2O8S2 * 6H2O}~;
    \item $M=\qty{392.13}{\g\per\mol}$~;
    \item P.F.~: $\qtyrange{100}{110}{\degreeCelsius}$~;
    \item \ce{(NH4)2Fe(SO4)2 * 6H2O}.
\end{itemize}
\begin{enumerate}
    \item Déduire de la formule de ce composé que l'élément fer s'y trouve sous
          forme d'ions \ce{Fe^{2+}}.
    \item Écrire l'équation de la réaction de dissolution du sel de
          \textsc{Mohr} dans l'eau.
    \item On désire préparer $\qty{100}{\mL}$ d'une solution telle que
          $\left[\ce{Fe^{2+}(aq)}\right]=\qty{0.020}{\mol\per\L}$. Quelle masse
          de sel de \textsc{Mohr} doit-on mettre en solution~?
\end{enumerate}

